\section{Motivation}\label{sec:introduction:motivation}
Existing work on quantum programming language design recognises a three-layer model separating the programming layer, the execution layer and the hardware layer~\cite{DiMatteo2024AbstractionHierarchy, nunez_corrales2025betterabstractmachines}.
In mature classical systems, these layers are separated by stable abstraction boundaries.
A programmer querying a relational database need not reason about B-trees or buffer pools to write correct queries.
Likewise, the programming layer of a quantum system should not require knowledge of the execution and hardware layers below it.
That separation has not been achieved in quantum programming.
Concepts from the execution and hardware layers consistently bleed through into the programming surface.
Programmers must select gate sequences, manage qubit allocation, reason about circuit depth and schedule uncomputation~\cite{DiMatteo2024AbstractionHierarchy}.
The full treatment of the quantum programming language abstraction landscape is presented in Chapter~\ref{ch:literature_review}.

Researchers and practitioners have observed that quantum programming systems consistently impose a vocabulary of quantum-mechanical concepts on the programmer, even when the problem domain has no intrinsic quantum structure~\cite{Gay2006QPLSurvey, DiMatteo2024AbstractionHierarchy, ferreira2024qplusage}.
Programmers must declare qubits, compose gate sequences, reason about circuit structure and manage entanglement.
These obligations hold regardless of whether the algorithm being implemented requires quantum-mechanical reasoning at the level of abstraction where the programmer is working.
This thesis formalises that burden as \textit{quantum bias}: the degree to which a language's surface model requires the programmer to understand quantum-mechanical vocabulary to write correct programs.
Quantum bias is driven by bleed-through, the leakage of execution-layer and hardware-layer concerns into the programming surface~\cite{DiMatteo2024AbstractionHierarchy}.
A language with zero quantum bias is \textit{quantum-implicit}, a concept introduced and defined in Section~\ref{sec:introduction:research_questions}.
Chapter~\ref{ch:literature_review} presents a comparative assessment of the existing quantum programming language landscape through which the concept is applied.

The most recent high-level quantum programming languages have made targeted progress in reducing the vocabulary burden.
\textbf{Silq}~\cite{bichsel2020silq} automates a class of uncomputation through its type system, relieving the programmer of manually constructing adjoint gate sequences in supported cases.
\textbf{Quff}~\cite{wright2024quff} introduces runtime function inversion, removing the need to write the inverse of a function explicitly.
\textbf{Qutes}~\cite{faro2025qutes, reversiblelifetime2026} introduces scope-bounded uncomputation semantics and a \texttt{ref} keyword that approximates quantum-agnostic pass-by-reference.
Each represents genuine semantic progress at a specific, principled point.
Each also retains quantum-mechanical obligations in the general case: Silq requires\texttt{qfree}, \texttt{const} and \texttt{lifted} annotations; Quff requires explicit bit-width specification and phase notation; Qutes requires explicit gate instructions for algorithms outside its built-in set.

\textbf{Qiskit}~\cite{aleksandrowicz2019qiskit}, the most widely deployed quantum programming system, makes no such attempt.
83\% of its active programmers are experienced classical engineers for whom quantum vocabulary is the primary reported barrier to effective use~\cite{ferreira2024qplusage}.
No existing quantum programming system achieves a surface model agnostic of the underlying quantum substrate.